\section{Related Work}

Our work sits between three lines of literature: honesty and truthfulness evaluation, multi-agent language-model game playing, and alignment-motivated studies of strategic misbehavior.

Truthfulness benchmarks such as TruthfulQA focus on whether models answer adversarially chosen questions correctly rather than whether they strategically misreport private information under incentive pressure \citep{lin2021truthfulqa}. Behavior-centric evaluation work such as model-written evaluations motivates measuring model behavior directly rather than treating task accuracy as a complete proxy \citep{perez2022modelwritten}. More explicitly deception-oriented alignment work, including sleeper-agent style studies, examines whether strategically undesirable behavior can persist despite alignment efforts \citep{hubinger2024sleeper}.

Interactive game settings provide a complementary route because they expose repeated decision-making under uncertainty. Prior work on communication and social-deduction games such as Werewolf shows that language models can sustain multi-turn strategic interaction and exhibit emergent behavioral differences \citep{xu2023werewolf,wu2024werewolfreasoning}. Our setting differs in a way that matters methodologically: Bullshit yields directly verifiable lie and truth labels from hidden cards rather than requiring indirect inference from persuasion, voting, or latent-role reasoning.

The benchmark also connects to broader work on oversight, strategic communication, and instruction-following under pressure \citep{irving2018debate,kenton2024weakjudges,ouyang2022instructgpt}. Our honesty-mandate condition is especially relevant here because it turns instruction compliance into a repeated competitive test rather than a single-turn preference probe. The result is a benchmark that is narrower than a general theory of deception, but sharper as a controlled evaluation artifact.
